% =============================================================================
%  Thesis Sections — 3.5.5 & 4.8
%  Simulator control strategies: conception + evaluation
% =============================================================================
%  Prerequisites in preamble:
%    \usepackage{graphicx, booktabs, multirow, amsmath, hyperref}
%    \graphicspath{{ml-backend/simulator/}}
% =============================================================================


% ─────────────────────────────────────────────────────────────────────────────
% 3.5.5  Conception
% ─────────────────────────────────────────────────────────────────────────────

\subsection{Stratégies de contrôle du simulateur}
\label{sec:control-strategies-design}

The simulator is not just a way to replay what happened---it is a tool for
deciding what to do next.  To support that, it includes a catalog of nine
interventions that a decision-maker (or an automated agent) can activate
during a simulation run.  These range from launching a targeted donor
campaign, deploying mobile collection units and extending centre operating
hours on the supply side, to fast-tracking lab workflows, adding surge staff,
enabling emergency inter-hospital sharing and activating clinical conservation
protocols on the operations side.

Each action has a defined operational cost, an activation delay, a ramp-up
period and a set of effects on simulation parameters such as donor arrival
rate, lab processing speed or demand management.
Table~\ref{tab:action-catalog} lists the full catalog.

\begin{table}[htbp]
\centering
\caption{Action catalog.  \emph{Cost} is per activation;
         \emph{Delay} is the time before the action begins to take effect;
         \emph{Ramp} is the time to reach full intensity;
         \emph{Active} is total duration including ramp-up.}
\label{tab:action-catalog}
\small
\begin{tabular}{@{}llcccp{5.2cm}@{}}
\toprule
\textbf{Action} & \textbf{Cost} & \textbf{Delay} & \textbf{Ramp}
                 & \textbf{Active} & \textbf{Key effects} \\
\midrule
Donor campaign         & 60  & 12\,h  & 120\,h & 120\,h
    & Arrivals $+50\%$, show-up $\times1.20$, committed donors $42\%$ \\
Mobile unit            & 55  & 24\,h  & 96\,h  & 120\,h
    & Arrivals $+80\%$, 2 extra collection sites \\
Extend hours           & 35  & 12\,h  & 72\,h  & 72\,h
    & Arrivals $+88\%$, $+3$\,h extension per centre \\
Lab fast-track         & 75  & 6\,h   & 48\,h  & 72\,h
    & Lab speed $\times0.30$, release delay $\times0.18$ \\
Surge staff            & 70  & 12\,h  & 48\,h  & 72\,h
    & $+4$ nurses, $+2$ lab, $+2$ processing staff \\
Emergency sharing      & 65  & 24\,h  & 96\,h  & 60\,h
    & Cross-hospital sharing, reserve$\to$0, replenishment $\times8$ \\
Clinical conservation  & 55  & 6\,h   & 36\,h  & 48\,h
    & Effective demand $-72\%$ \\
National mutual aid    & 85  & 48\,h  & 120\,h & 120\,h
    & Transport relief $\times0.80$, replenishment $\times1.80$ \\
Rapid courier          & 48  & 2\,h   & 12\,h  & 36\,h
    & Transport relief $\times0.35$, replenishment $\times1.75$ \\
\bottomrule
\end{tabular}
\end{table}

These actions are grouped into \emph{strategies}.  The system provides five
static strategies that pre-activate fixed combinations of actions, and three
adaptive strategies driven by reinforcement-learning agents
(Table~\ref{tab:strategies}).

\begin{table}[htbp]
\centering
\caption{Built-in strategies.  Static strategies activate a fixed set of
         actions at the start of the episode.  Adaptive strategies let a
         trained agent choose which actions to activate and at what intensity
         every 6 simulated hours.}
\label{tab:strategies}
\small
\begin{tabular}{@{}llp{6.8cm}@{}}
\toprule
\textbf{Strategy} & \textbf{Type} & \textbf{Actions / controller} \\
\midrule
Baseline               & Static   & No intervention \\
Mass donor campaign    & Static   & Campaign + mobile unit + extend hours \\
Lab investment         & Static   & Lab fast-track \\
Emergency network      & Static   & Surge staff + emergency sharing
                                    + clinical conservation \\
Full crisis response   & Static   & All 8 actions simultaneously \\
\addlinespace
PPO Shortage Minimizer & Adaptive & Discrete PPO; chooses 1 of 10 options
                                    (9~actions + no-op) \\
PPO Continuous         & Adaptive & PPO in $\text{Box}(-1,1)^{14}$;
                                    controls all intensities \\
Official DreamerV3     & Adaptive & Model-based RL; same continuous action
                                    space as PPO Continuous \\
\bottomrule
\end{tabular}
\end{table}

The adaptive strategies observe the simulation state every six simulated
hours through a 16-dimensional vector (Table~\ref{tab:obs-space}).  Based on
this observation the agent decides which actions to activate and at what
intensity.

\begin{table}[htbp]
\centering
\caption{Observation vector ($\text{Box}(0,1)^{16}$).  All values are
         clipped to $[0,1]$.}
\label{tab:obs-space}
\small
\begin{tabular}{@{}clp{6.4cm}@{}}
\toprule
\textbf{Index} & \textbf{Feature} & \textbf{Description} \\
\midrule
0  & recent shortage rate     & Shortage / requested in current step \\
1  & cumulative shortage rate & Total shortage / total net requested \\
2  & critical gap             & Worst component-coverage gap vs.\ target \\
3  & RBC gap                  & (Target $-$ RBC coverage days) / target \\
4  & platelet gap             & Same, for platelets ($\times1.15$ weight) \\
5  & plasma gap               & Same, for plasma ($\times0.65$ weight) \\
6  & transport stress         & Effective transport penalty normalised \\
7  & weather stress           & 0\,(clear) to 1\,(blizzard) \\
8  & donor stress             & Inter-arrival and show-up deficit \\
9  & demand stress            & Demand rate and surge factor excess \\
10 & service pressure         & Recent requested / transfused \\
11 & active cost ratio        & Action cost / 500 \\
12 & forecast signal          & Demand forecast pressure \\
13 & priority pressure        & Priority-weighted shortage signal \\
14 & congestion stress        & Transport congestion level \\
15 & budget pressure          & Fraction of budget consumed \\
\bottomrule
\end{tabular}
\end{table}

\paragraph{Choice of agent architecture.}

The question is what kind of agent is best suited for this control problem.
We considered three architectures.

The first is a \textbf{PPO agent with discrete actions}.  At each decision
step it picks one of ten options (nine actions plus a no-op).  This is
simple to train but limited: it can only activate one action at a time,
it cannot control intensity, and it has no internal model of what will
happen next.  It reacts to the current state but does not plan.

The second is \textbf{PPO with continuous actions}.  This version operates
in $\text{Box}(-1,1)^{14}$ where each dimension controls the intensity of
one intervention, one hospital routing dial or one component allocation
dial.  It can activate multiple actions simultaneously and fine-tune their
levels, which is a better fit for the problem.  But it still has no world
model---it maps observations to actions without modelling how the supply
chain will evolve.

The third is \textbf{DreamerV3}~\cite{hafner2023dreamerv3}, a model-based
reinforcement-learning agent.  DreamerV3 learns a latent world model from
experience: it builds an internal representation of how the environment
transitions from one state to another, and it uses that model to
\emph{imagine} future trajectories before deciding what to do.  This matters
for blood supply management because the effects of interventions are
delayed---a donor campaign takes 12~hours to activate and 5~days to ramp
up---so an agent that can plan ahead should outperform one that only reacts
to what it sees right now.  DreamerV3 operates in the same continuous action
space as PPO~Continuous.

\paragraph{Reward function.}

The reward function that guides all three agents is structured in three
tiers to reflect operational priorities
(Table~\ref{tab:reward-tiers}).

\begin{table}[htbp]
\centering
\caption{Per-step reward tiers (DreamerV3 variant).  The raw reward is
         centred at $-0.55$ and squashed through
         $\tanh\!\bigl((r-0.55)\times2.5\bigr)$, mapping a perfect step
         to ${\approx}+0.95$ and a catastrophic step to
         ${\approx}-0.99$.}
\label{tab:reward-tiers}
\small
\begin{tabular}{@{}clcl@{}}
\toprule
\textbf{Tier} & \textbf{Focus} & \textbf{Weight} & \textbf{Key terms} \\
\midrule
1 & Patient safety      & ${\sim}55\%$
  & $+1.00\times\text{fill\_rate}$, \;
    $-0.60\times\text{priority\_shortage}$, \;
    $-0.40\times\text{incompatible}$ \\
2 & Matching quality    & ${\sim}20\%$
  & $+0.25\times\text{exact\_match}$, \;
    $+0.05\times\text{compatible\_sub}$ \\
3 & Waste \& cost       & ${\sim}25\%$
  & $-0.35\times\text{wastage}$, \;
    $-0.08\times\text{budget\_overspend}$, \;
    $-0.07\times\text{congestion}$ \\
\bottomrule
\end{tabular}
\end{table}

This structure ensures the agent prioritises keeping patients supplied, then
optimises for transfusion quality, then manages costs---in that order.  The
choice of DreamerV3 as the primary adaptive strategy was motivated by three
factors: the delayed-effect nature of the action space, the need for
simultaneous multi-action control, and the expected benefit of internal
planning over pure reaction.  Whether this design reasoning holds up against
the data is evaluated in Section~\ref{sec:strategy-eval}.


% ─────────────────────────────────────────────────────────────────────────────
% 4.8  Evaluation
% ─────────────────────────────────────────────────────────────────────────────

\section{Évaluation des stratégies de contrôle}
\label{sec:strategy-eval}

\subsection{Protocole d'évaluation}
\label{sec:eval-protocol}

Before comparing strategies we first validated that the simulator itself
produces realistic outputs.  A 30-replication experiment on the baseline
scenario (seeds~1--30, 336 simulated hours each) was compared against
published Héma-Québec and Statistics~Canada benchmarks.
Table~\ref{tab:validation-results} summarises the results.

\begin{table}[htbp]
\centering
\caption{Simulator validation against real-world references (30
         replications, baseline scenario).  CI~=~95\% confidence interval.
         MAPE and Theil's~U are computed against the reference midpoint.}
\label{tab:validation-results}
\small
\begin{tabular}{@{}lccccc@{}}
\toprule
\textbf{Metric}
    & \textbf{Sim.\ mean}
    & \textbf{95\% CI}
    & \textbf{Reference}
    & \textbf{MAPE}
    & \textbf{Theil U} \\
\midrule
Daily donations          & 39.06 & [38.06,\;40.07] & 39.69  &  5.8\% & 0.035 \\
Shortage rate (\%)       &  7.41 & [5.61,\;9.22]   &  5.00  & 80.7\% & 0.386 \\
Deferral rate (\%)       & 12.63 & [11.82,\;13.44]  & 12.00  & 16.3\% & 0.090 \\
Wastage rate (\%)        &  4.46 & [3.62,\;5.30]   &  3.00  & 68.9\% & 0.332 \\
Service fulfilment (\%)  & 92.59 & [90.78,\;94.39]  & 97.00  &  5.1\% & 0.034 \\
Exact match rate (\%)    & 75.44 & [74.25,\;76.64]  & 80.00  &  5.9\% & 0.036 \\
\bottomrule
\end{tabular}
\end{table}

Figure~\ref{fig:validation-bar} shows the simulated values normalised
against reference midpoints.  Daily completed donations land at
$0.99\times$ the reference, deferral rate at $1.15\times$, service
fulfilment at $0.96\times$.  The largest deviation is shortage rate at
$1.41\times$ (7.41\% vs.\ 5.0\%), though the simulated 95\% CI still
overlaps the real-world range of 0.5--10.0\%.
Figure~\ref{fig:ci-coverage} confirms that for all six metrics the
simulation CI falls within or closely overlaps the published reference
range.  Blood-type proportions match the Canadian population distribution
with a chi-squared $p$-value of 0.923.  All 18 validation tests pass.

\begin{figure}[htbp]
\centering
\includegraphics[width=0.78\linewidth]
    {validation_results/plots/validation_comparison_bar.png}
\caption{Simulated outputs normalised against real-world reference
         midpoints (dashed line~$=$~1.0).  Blue~=~simulated,
         grey~=~reference.}
\label{fig:validation-bar}
\end{figure}

\begin{figure}[htbp]
\centering
\includegraphics[width=0.78\linewidth]
    {validation_results/plots/ci_coverage_plot.png}
\caption{95\% confidence interval coverage.
         Green~=~simulation CI, grey~=~real-world reference range,
         diamonds~=~point estimates.}
\label{fig:ci-coverage}
\end{figure}

A sensitivity analysis (Figure~\ref{fig:tornado}) confirmed that the three
most influential parameters are demand-side:
\texttt{demand\_surge\_factor} swings shortage rate by 25 percentage points
when varied $\pm20\%$, followed by \texttt{avg\_units\_per\_order}
(23\,pp) and \texttt{demand\_rate\_h} (20\,pp).  Supply-side parameters
like \texttt{eligible\_rate} and \texttt{transport\_penalty} have less than
2\,pp of impact.  This tells us that any effective strategy must primarily
manage demand pressure---which is exactly what the adaptive agents are
designed to do.

\begin{figure}[htbp]
\centering
\includegraphics[width=0.78\linewidth]
    {validation_results/plots/tornado_diagram.png}
\caption{Tornado diagram---sensitivity of shortage rate to $\pm20\%$
         parameter perturbation.  Blue~=~parameter decreased,
         red~=~parameter increased.}
\label{fig:tornado}
\end{figure}

For the strategy comparison itself, every combination of
$9$~scenarios~$\times$~$5$~strategies was run with seed~42 and a 168-hour
horizon.  The five strategies compared are: baseline (no intervention),
PPO~Shortage Minimizer (discrete actions), PPO~Continuous (continuous
actions), Official~DreamerV3 (model-based, continuous) and DreamerV4
(early-stage checkpoint, included for completeness).  The static
strategies (mass campaign, lab investment, emergency network, full crisis
response) were evaluated separately in an extended 7-strategy comparison.


% ─────────────────────────────────────────────────────────────────────────────
\subsection{Résultats comparatifs}
\label{sec:eval-results}

Figure~\ref{fig:heatmaps} shows the performance heatmaps across all
scenarios.  The results are clear: \textbf{DreamerV3 achieves the lowest
shortage rate in every single scenario}, often by a wide margin.

\begin{figure}[htbp]
\centering
\includegraphics[width=\linewidth]
    {results/strategy_eval_heatmaps.png}
\caption{Strategy performance heatmaps---5 adaptive/baseline strategies
         across 9 scenarios (seed~42, 168\,h).  Darker colours indicate
         worse values for shortage-related metrics and lower costs for
         budget metrics.}
\label{fig:heatmaps}
\end{figure}

Table~\ref{tab:shortage-comparison} extracts the shortage rates for a
representative subset of scenarios.

\begin{table}[htbp]
\centering
\caption{Shortage rate (\%) across selected scenarios.  Best result per
         scenario in \textbf{bold}.  Training scenarios above the line,
         holdout scenarios below.}
\label{tab:shortage-comparison}
\small
\begin{tabular}{@{}lccccc@{}}
\toprule
\textbf{Scenario}
    & \textbf{Baseline}
    & \textbf{PPO\,SM}
    & \textbf{PPO\,Cont.}
    & \textbf{DreamerV3}
    & \textbf{DreamerV4} \\
\midrule
Normal operations        &  5.0 &  4.4 &  3.9 & \textbf{0.5}  &  6.6 \\
Donor decline            & 22.9 &  2.0 &  5.0 & \textbf{0.3}  & 22.7 \\
Demand surge             & 60.2 & 19.3 & 27.0 & \textbf{9.4}  & 53.4 \\
Transport disruption     & 65.4 & 22.6 &  7.1 & \textbf{2.2}  & 59.9 \\
Combined crisis          & 80.5 & 55.6 & 40.2 & \textbf{23.0} & 78.9 \\
\midrule
Holiday flu wave         & 34.2 &  7.1 &  5.7 & \textbf{2.9}  & 30.0 \\
Regional testing backlog &  9.4 &  3.3 &  4.1 & \textbf{0.5}  &  7.6 \\
Provincial supply crunch & 27.5 &  5.9 &  5.6 & \textbf{0.8}  & 17.4 \\
Post-storm backlog       & 44.2 &  4.5 &  7.4 & \textbf{1.5}  & 36.3 \\
\bottomrule
\end{tabular}
\end{table}

On the baseline scenario DreamerV3 reaches 0.5\% shortage versus 5.0\% for
the passive baseline and 3.9\% for PPO~Continuous.  On the most severe
combined crisis, DreamerV3 keeps shortage at 23.0\% where the passive
baseline hits 80.5\% and PPO~Continuous reaches 40.2\%.

The holdout scenarios---which mix stressors in ways the agents never saw
during training---tell the same story.  On the holiday flu wave DreamerV3
achieves 2.9\% versus 34.2\% for baseline.  On the provincial supply crunch
0.8\% versus 27.5\%.  On the post-storm surgical backlog 1.5\% versus
44.2\%.

DreamerV3 also leads on exact-match transfusion rate: 84.2\% on normal
operations, 85.8\% on donor decline, 87.5\% on regional testing backlog---
consistently above the other strategies.  It does not just prevent
shortages; it fulfils orders with better-matched blood.

The extended comparison including static strategies
(Figure~\ref{fig:all-heatmaps}) reveals an important pattern.  The full
crisis response strategy---which activates all eight interventions at
once---achieves strong shortage reduction (0.1\% on baseline, 27.0\% on
donor decline) but at the cost of massive expiration: 182~expired units on
baseline, 170 on donor decline, 224 on provincial supply crunch.
DreamerV3 reaches comparable or better shortage rates with fewer expired
units in crisis scenarios, while also adapting its interventions dynamically
rather than spending every resource at once.

\begin{figure}[htbp]
\centering
\includegraphics[width=\linewidth]
    {results/strategy_eval_all_heatmaps.png}
\caption{Extended heatmaps---7 strategies (5 static $+$ 2 adaptive) across
         9~scenarios.  The full crisis response strategy reduces shortage
         aggressively but causes high unit expiration.}
\label{fig:all-heatmaps}
\end{figure}

The tradeoff between shortage reduction and wastage is visible on the
Pareto frontier (Figure~\ref{fig:pareto}).  DreamerV3 sits at the far-left
edge---lowest average shortage rate (${\sim}5\%$)---but with higher total
expired units (${\sim}55$) than baseline (${\sim}10$) or
PPO~Shortage~Minimizer (${\sim}12$).  This is the cost of aggressive
stocking: some units expire before they are needed.  PPO~Continuous occupies
a middle ground (${\sim}22\%$ shortage, ${\sim}25$ expired).  DreamerV4 is
dominated---high shortage, high action cost and no advantage on any axis.

\begin{figure}[htbp]
\centering
\includegraphics[width=0.82\linewidth]
    {results/strategy_eval_pareto_overall.png}
\caption{3-D Pareto frontier across all scenarios.  Axes: shortage rate,
         action cost, total expired (all lower-is-better).  Circled points
         are Pareto-optimal.}
\label{fig:pareto}
\end{figure}

\paragraph{Budget sensitivity.}

The budget sensitivity analysis (Figure~\ref{fig:budget-lines}) shows how
strategies degrade as budgets tighten.  DreamerV3 is the most
budget-responsive agent: at a budget of 12\,000 it drives shortage below
5\%; at 7\,500 shortage rises to about 7\%; at 3\,000 it climbs to about
38\%.  PPO~Shortage~Minimizer is more stable but plateaus higher---around
20\% shortage regardless of budget above 6\,000.  Baseline sits flat at
${\sim}60\%$ since it never spends anything.

The efficiency frontier (Figure~\ref{fig:efficiency}) makes this clearer:
for any given level of budget spent, DreamerV3 achieves a lower shortage
rate than every other strategy.  At a spend of ${\sim}5\,000$, DreamerV3
reaches 5\% shortage while PPO~Shortage~Minimizer sits at 20\% and
PPO~Continuous at 25\%.

\begin{figure}[htbp]
\centering
\includegraphics[width=0.82\linewidth]
    {results/budget_eval_budget_lines.png}
\caption{Strategy performance as a function of episode budget (averaged
         across scenarios).  Shaded bands show the spread across scenarios.
         DreamerV3 (orange) improves steeply with additional budget while
         Baseline (blue) is unaffected.}
\label{fig:budget-lines}
\end{figure}

\begin{figure}[htbp]
\centering
\includegraphics[width=0.82\linewidth]
    {results/budget_eval_efficiency_frontier.png}
\caption{Budget efficiency frontier.  Each point is one
         (strategy,\,budget-level) pair.  DreamerV3 dominates the lower
         envelope---for any amount spent it achieves the lowest shortage
         rate.}
\label{fig:efficiency}
\end{figure}

\paragraph{Training convergence.}

The DreamerV3 training curves (Figure~\ref{fig:training}) confirm that the
agent learned effectively.  The \texttt{m1\_continuous\_run8\_kpi\_aligned}
run, trained for 1~million environment steps (${\sim}10\,000$ episodes),
shows episode scores rising from near-zero to a stable average of 67.4 in
the final training phase, with peaks above 100.  The value loss and
continuation loss both converge steadily, indicating the world model
learned a useful internal representation of the blood supply chain
dynamics.  The reward rate stabilises around 0.80 in the second half of
training.

\begin{figure}[htbp]
\centering
\includegraphics[width=0.82\linewidth]
    {dreamerv3_runs/m1_continuous_run8_kpi_aligned/metrics_plot.png}
\caption{DreamerV3 training metrics over 1\,M steps.
         Top-left: episode score; top-right: reward rate;
         bottom-left: value loss; bottom-right: continuation loss.
         Both losses converge, confirming the world model learned a
         useful latent representation.}
\label{fig:training}
\end{figure}


% ─────────────────────────────────────────────────────────────────────────────
\subsection{Justification du choix de DreamerV3}
\label{sec:dreamerv3-justification}

The evaluation results confirm the design hypothesis from
Section~\ref{sec:control-strategies-design}.  DreamerV3 outperforms every
other strategy on shortage rate across all nine scenarios, including the
four holdout scenarios it was never explicitly trained on.  Three properties
explain this.

\paragraph{Planning through a world model.}

Blood supply interventions have activation delays ranging from 2 to
48~hours and ramp-up periods of 12 to 120~hours.  A purely reactive agent
(PPO) sees the shortage happening and fires an action, but by the time the
action takes effect the crisis may have already deepened.  DreamerV3
imagines future trajectories in its learned world model and acts
pre-emptively---which is why its shortage rates are consistently lower even
in fast-onset scenarios like the demand surge (9.4\% vs.\ 27.0\% for
PPO~Continuous).

\paragraph{Simultaneous multi-action control.}

The discrete PPO agent can only activate one action per decision step.
DreamerV3 controls all nine intervention intensities, two hospital routing
dials and three component allocation dials simultaneously through a
continuous action vector.  The heatmaps show this matters: DreamerV3
manages to reduce net requested units (through demand management) while
also boosting supply (through campaigns and replenishment)---something a
single-action controller cannot coordinate.  For instance, on the transport
disruption scenario DreamerV3 serves 738~net-requested units with a 2.2\%
shortage, while PPO~Shortage~Minimizer faces 1\,179~units with a 22.6\%
shortage because it cannot activate conservation and replenishment at the
same time.

\paragraph{Generalisation to unseen crises.}

The holdout scenarios were designed to combine stressors in novel ways.
DreamerV3's world model lets it transfer what it learned about individual
stressors (donor decline, transport disruption, demand pressure) to new
combinations it has never encountered.  On the post-storm surgical
backlog---a scenario that pairs rebounding demand with slow donor recovery
and damaged transport---DreamerV3 achieves 1.5\% shortage where the
next-best adaptive strategy (PPO~Shortage~Minimizer) reaches 4.5\% and
baseline hits 44.2\%.

\paragraph{Tradeoff: wastage.}

The main tradeoff is wastage.  DreamerV3's aggressive stocking leads to
more expired units than conservative strategies
(Table~\ref{tab:tradeoff}).  This is a conscious design choice embedded in
the reward function: Tier~1 (patient safety) carries ${\sim}55\%$ of the
reward weight while Tier~3 (waste and cost) carries ${\sim}25\%$.  In a
blood supply context a shortage means a patient does not receive the
transfusion they need, while an expired unit is an economic loss.  The
reward structure reflects the judgement that the former is worse than the
latter.

\begin{table}[htbp]
\centering
\caption{Shortage--wastage tradeoff (average across 9 scenarios).}
\label{tab:tradeoff}
\small
\begin{tabular}{@{}lccc@{}}
\toprule
\textbf{Strategy}
    & \textbf{Shortage (\%)}
    & \textbf{Expired units}
    & \textbf{Action cost} \\
\midrule
Baseline               & 38.8 &  6.6 &   0.0 \\
PPO Shortage Minimizer & 13.9 &  7.3 & 177.4 \\
PPO Continuous         & 11.7 & 10.6 &  50.2 \\
\textbf{DreamerV3}     & \textbf{4.6}  & 34.8 & 158.2 \\
DreamerV4              & 36.3 & 12.6 & 198.2 \\
\bottomrule
\end{tabular}
\end{table}
